\documentclass[aspectratio=169,11pt]{beamer}
% Shared Beamer style for practice contest solution walkthroughs.
\usepackage[utf8]{inputenc}
\usepackage[T1]{fontenc}
\usepackage{lmodern}
\usepackage{amsmath,amssymb}
\usepackage{xcolor}
\usepackage{hyperref}
\usepackage{listings}

\definecolor{CPNavy}{HTML}{172033}
\definecolor{CPBlue}{HTML}{2563EB}
\definecolor{CPGreen}{HTML}{16A34A}
\definecolor{CPOrange}{HTML}{F97316}
\definecolor{CPGray}{HTML}{64748B}
\definecolor{CPLight}{HTML}{F8FAFC}
\definecolor{CPCodeBg}{HTML}{F1F5F9}
\definecolor{CPCodeRule}{HTML}{CBD5E1}

\usetheme{default}
\usefonttheme{professionalfonts}
\setbeamertemplate{navigation symbols}{}
\setbeamersize{text margin left=0.65cm,text margin right=0.65cm}
\setbeamertemplate{itemize item}{\color{CPBlue}$\blacktriangleright$}
\setbeamertemplate{itemize subitem}{\color{CPGray}$\triangleright$}

\setbeamercolor{normal text}{fg=CPNavy,bg=white}
\setbeamercolor{frametitle}{fg=white,bg=CPNavy}
\setbeamercolor{title}{fg=white,bg=CPNavy}
\setbeamercolor{subtitle}{fg=CPLight,bg=CPNavy}
\hypersetup{colorlinks=true,urlcolor=CPBlue}

\setbeamertemplate{frametitle}{%
  \nointerlineskip%
  \begin{beamercolorbox}[wd=\paperwidth,ht=0.88cm,dp=0.22cm,leftskip=0.55cm,rightskip=0.35cm]{frametitle}
    \usebeamerfont{frametitle}\insertframetitle\hfill\small\insertframenumber/\inserttotalframenumber
  \end{beamercolorbox}%
}

\setbeamertemplate{footline}{%
  \leavevmode%
  \hbox{%
    \begin{beamercolorbox}[wd=.58\paperwidth,ht=2.4ex,dp=1ex,leftskip=.5cm]{author in head/foot}%
      \color{CPGray}\insertshorttitle
    \end{beamercolorbox}%
    \begin{beamercolorbox}[wd=.42\paperwidth,ht=2.4ex,dp=1ex,rightskip=.5cm plus1fil]{date in head/foot}%
      \hfill\color{CPGray}\insertshortauthor
    \end{beamercolorbox}%
  }%
}

\setbeamertemplate{title page}{%
  \vspace*{1.25cm}
  \begin{beamercolorbox}[wd=\paperwidth,sep=0.65cm,leftskip=0.15cm,rightskip=0.15cm]{title}
    {\color{white}\Huge\bfseries\inserttitle\par}
    \vspace{0.35cm}
    {\color{CPLight}\Large\insertsubtitle\par}
    \vspace{0.6cm}
    {\color{CPLight}\large\insertauthor\par}
  \end{beamercolorbox}
  \vfill
}

\newcommand{\sectiontag}[1]{\textbf{\color{CPBlue}#1}}
\newenvironment{tightitemize}{\begin{itemize}\small\setlength{\itemsep}{0.18cm}}{\end{itemize}}
\lstdefinestyle{cpcode}{
  language=Python,
  basicstyle=\ttfamily\tiny,
  keywordstyle=\color{CPBlue}\bfseries,
  commentstyle=\color{CPGray}\itshape,
  stringstyle=\color{CPGreen},
  backgroundcolor=\color{CPCodeBg},
  frame=single,
  framerule=0.25pt,
  rulecolor=\color{CPCodeRule},
  showstringspaces=false,
  columns=fullflexible,
  keepspaces=true,
  breaklines=true,
  breakatwhitespace=false,
  tabsize=4,
  xleftmargin=0.08cm,
  xrightmargin=0.08cm,
  aboveskip=0.08cm,
  belowskip=0.08cm
}


\title[shortestpath1]{Single Source Shortest Path}
\subtitle{Day 4 solution walkthrough}
\author[Solution Deck]{Competitive Programming Bootcamp}
\date{}

\begin{document}

\begin{frame}[plain]
  \titlepage
\end{frame}

% BEGIN GENERATED STATEMENT INTERPRETATION
\begin{frame}{Interpret the Word Problem}
\small
\sectiontag{Statement excerpts:}
\begin{tightitemize}
  \item \textit{``minimum distance from node \$s\$''}
  \item \textit{``if there is no path from \$s\$ to that node''}
\end{tightitemize}

\vspace{0.18cm}
\sectiontag{Programming interpretation:}
\begin{tightitemize}
  \item Each test case is a directed graph with nonnegative edge weights.
  \item Run Dijkstra once from the source, then answer all destination queries from dist.
  \item The all-zero line ends the input.
\end{tightitemize}
\end{frame}
% END GENERATED STATEMENT INTERPRETATION







\begin{frame}{Problem Model}
\small
\sectiontag{Textbook connection:} CP4 4.4.3 Dijkstra's Algorithm

\vspace{0.25cm}
\sectiontag{Core technique:} Standard Dijkstra with query answering

\vspace{0.25cm}
\begin{tightitemize}
  \item Each test case is a directed weighted graph.
  \item All edge weights are nonnegative.
  \item After one source is fixed, multiple queries ask for distances from that source.
\end{tightitemize}
\end{frame}

\begin{frame}{How to Recognize the Approach}
\begin{tightitemize}
  \item Nonnegative weighted shortest paths point to Dijkstra.
  \item Multiple queries from the same source should be answered after one run.
  \item An adjacency list stores only the given directed edges.
\end{tightitemize}
\end{frame}

\begin{frame}{Algorithm}
\begin{tightitemize}
  \item Read V, E, q, and source s until 0 0 0 0.
  \item Build AL[u] as a list of (v, w) directed edges.
  \item Initialize dist[s] = 0 and all others to INF.
  \item Use a priority queue to process vertices by smallest known distance.
  \item Relax outgoing edges and answer each query from the final dist array.
\end{tightitemize}
\end{frame}

\begin{frame}{Walkthrough of the Key State}
\begin{tightitemize}
  \item A stale priority queue entry has d > dist[u] and can be skipped.
  \item Unreachable query vertices remain at INF.
  \item The output requires a blank line after each test case.
\end{tightitemize}
\end{frame}

\begin{frame}{Why the Algorithm Is Correct}
\begin{tightitemize}
  \item Dijkstra's invariant finalizes each popped non-stale vertex at its true shortest distance.
  \item Every directed edge is relaxed when its source is finalized or improved.
  \item The distance array therefore answers all post-run queries correctly.
\end{tightitemize}
\end{frame}

\begin{frame}{Complexity and Implementation Notes}
\small
\sectiontag{Complexity:} O((V + E) log V) time per test case and O(V + E) memory.

\vspace{0.3cm}
\begin{tightitemize}
  \item Do not add reverse edges; the graph is directed.
  \item Use a large integer INF instead of float infinity for integer distances.
\end{tightitemize}
\end{frame}



% BEGIN GENERATED CODE WALKTHROUGH
\begin{frame}[fragile]{Code: Read Test Case Header}
\scriptsize
\sectiontag{Source lines:} \texttt{shortestpath1.py:48--65}

\vspace{0.08cm}
\begin{columns}[T,totalwidth=\textwidth]
\begin{column}{0.58\textwidth}
\begin{lstlisting}[style=cpcode]
import sys
from heapq import heappush, heappop

def main():
    input = sys.stdin.buffer.readline

    INF = int(1e18)

    output = []

    while True:
        V, E, q, s = map(int, input().split())

        if V == 0 and E == 0 and q == 0 and s == 0:
            break
\end{lstlisting}
\end{column}
\begin{column}{0.38\textwidth}
\sectiontag{What this chunk does}
\begin{tightitemize}
  \item The loop continues until the sentinel 0 0 0 0.
  \item INF represents vertices that have not been reached.
  \item output buffers answers across test cases.
\end{tightitemize}
\end{column}
\end{columns}
\end{frame}

\begin{frame}[fragile]{Code: Build Directed Adjacency}
\scriptsize
\sectiontag{Source lines:} \texttt{shortestpath1.py:67--78}

\vspace{0.08cm}
\begin{columns}[T,totalwidth=\textwidth]
\begin{column}{0.58\textwidth}
\begin{lstlisting}[style=cpcode]
AL = [[] for u in range(V)]

for _ in range(E):
    u, v, w = map(int, input().split())

    AL[u].append((v, w))
\end{lstlisting}
\end{column}
\begin{column}{0.38\textwidth}
\sectiontag{What this chunk does}
\begin{tightitemize}
  \item AL[u] stores only edges leaving u.
  \item The edge direction matters; v to u is not added unless it appears separately.
  \item Each stored pair is destination and weight.
\end{tightitemize}
\end{column}
\end{columns}
\end{frame}

\begin{frame}[fragile]{Code: Initialize Dijkstra}
\scriptsize
\sectiontag{Source lines:} \texttt{shortestpath1.py:79--93}

\vspace{0.08cm}
\begin{columns}[T,totalwidth=\textwidth]
\begin{column}{0.58\textwidth}
\begin{lstlisting}[style=cpcode]
dist = [INF for u in range(V)]

dist[s] = 0

pq = []
heappush(pq, (0, s))
\end{lstlisting}
\end{column}
\begin{column}{0.38\textwidth}
\sectiontag{What this chunk does}
\begin{tightitemize}
  \item All distances start at INF except the source distance of zero.
  \item The heap orders vertices by best known distance.
  \item The initial heap entry starts the search from s.
\end{tightitemize}
\end{column}
\end{columns}
\end{frame}

\begin{frame}[fragile]{Code: Relax Edges}
\scriptsize
\sectiontag{Source lines:} \texttt{shortestpath1.py:95--122}

\vspace{0.08cm}
\begin{columns}[T,totalwidth=\textwidth]
\begin{column}{0.58\textwidth}
\begin{lstlisting}[style=cpcode]
while len(pq) > 0:
    d, u = heappop(pq)

    if d > dist[u]:
        continue

    for v, w in AL[u]:

        new_distance = dist[u] + w

        if new_distance >= dist[v]:
            continue

        dist[v] = new_distance

        heappush(pq, (dist[v], v))
\end{lstlisting}
\end{column}
\begin{column}{0.38\textwidth}
\sectiontag{What this chunk does}
\begin{tightitemize}
  \item The smallest queued distance is popped first.
  \item Stale heap entries are ignored.
  \item Each edge relaxation improves dist[v] and pushes the new candidate distance.
\end{tightitemize}
\end{column}
\end{columns}
\end{frame}

\begin{frame}[fragile]{Code: Answer Queries}
\scriptsize
\sectiontag{Source lines:} \texttt{shortestpath1.py:124--143}

\vspace{0.08cm}
\begin{columns}[T,totalwidth=\textwidth]
\begin{column}{0.58\textwidth}
\begin{lstlisting}[style=cpcode]
        for _ in range(q):
            vertex = int(input())

            if dist[vertex] == INF:
                output.append("Impossible")
            else:
                output.append(str(dist[vertex]))

        output.append("")

    sys.stdout.write("\n".join(output))

main()
\end{lstlisting}
\end{column}
\begin{column}{0.38\textwidth}
\sectiontag{What this chunk does}
\begin{tightitemize}
  \item After Dijkstra, each query is just one array lookup.
  \item Unreached vertices still have INF and print Impossible.
  \item A blank line is added after each test case as required.
\end{tightitemize}
\end{column}
\end{columns}
\end{frame}
% END GENERATED CODE WALKTHROUGH

\begin{frame}{Source}
\small
\sectiontag{Solution file:} \texttt{practice\_contests/day\_4/shortestpath1.py}

\vspace{0.3cm}
\sectiontag{Problem:} \url{https://open.kattis.com/problems/shortestpath1}

\vspace{0.45cm}
Use this deck with the source code open beside it. The slides explain the reasoning; the code shows the exact input handling and output formatting.
\end{frame}

\end{document}
